\documentclass[10pt,a4paper]{article}
\usepackage{amsmath,amssymb,amsthm,graphicx}
\usepackage{hyperref}
\usepackage[capitalize, nameinlink]{cleveref}
\usepackage{amssymb}
\usepackage[bbgreekl]{mathbbol}
\usepackage{graphicx}
\usepackage{enumitem}
\usepackage{xifthen}
\usepackage{comment}
\newcommand\myitem[1][]{\item[#1]\refstepcounter{dummy}\def\@currentlabel{#1}}
\newcommand{\thistheoremname}{}
\newtheorem*{genericthm*}{\thistheoremname}
\newenvironment{namedthm*}[1]
{\renewcommand{\thistheoremname}{#1}%
	\begin{genericthm*}}
	{\end{genericthm*}}

% Mathematical macros
\newcommand{\absolutevalue}[1]{|#1|}
\newcommand{\ADR}{\mathtt{AD}_{\R}}
\DeclareMathOperator{\appr}{appr}
\newcommand{\Bairespace}[1][]{
	\ifthenelse{\equal{#1}{}}{\functions{\N}{\N}}{\functions{#1}{\N}}
}
\newcommand{\bbL}{\mathbb{L}}
\newcommand{\bbX}{\mathbb{X}}
\newcommand{\Borelhomomorphism}{\le_B}
\newcommand{\Cantorspace}[1][]{
	\ifthenelse{\equal{#1}{}}{\functions{\N}{2}}{\functions{#1}{2}}
}
\newcommand{\completegraph}[1]{K_{#1}}
\newcommand{\constantsequence}[2]{\sequence{#1}^{#2}}
\newcommand{\continuoushomomorphism}{\le_c}
\newcommand{\definedterm}[1]{\emph{#1}}
\newcommand{\domain}[1]{\mathrm{dom}(#1)}
\newcommand{\from}{\colon}
\newcommand{\goesto}{\rightarrow}
\newcommand{\graphmetric}[1]{d_{#1}}
\newcommand{\Gzero}[1][]{
	\ifthenelse{\equal{#1}{}}
	{\mathbb{G}_0}
	{\mathbb{G}_{0,n}}
}
\newcommand{\Lzero}[1][]{\ifthenelse{\equal{#1}{}}{\bbL_0}{L_{0, #1}}}
\newcommand{\Lzerospace}[1][]{\ifthenelse{\equal{#1}{}}{\bbX_0}{X_{0, #1}}}
\newcommand{\mathand}{\text{ and }}
\newcommand{\modulo}[1]{\ (\text{mod } 2)}
\newcommand{\N}{\mathbb{N}}
\newcommand{\odometer}[1]{\mathbb{\Sigma}_{#1}}
\newcommand{\pair}[2]{(#1, #2)}
\newcommand{\partialto}{\rightharpoonup}
\newcommand{\pathlength}[1]{\mathrm{length}(#1)}

\newcommand{\rbar}{\bar{r}}
\newcommand{\rr}{\textcolor{red}}
\newcommand{\sbar}{\bar{s}}
\newcommand{\shift}[2]{\mathbb{S}_{#1, #2}}
\newcommand{\signedlength}[2]{\length{#1}_{#2}}
\newcommand{\suchthat}{\mid}
\newcommand{\R}{\mathbb{R}}
\renewcommand{\restriction}[2]{#1 \upharpoonright #2}

\newcommand{\triple}[3]{(#1, #2, #3)}
\newcommand{\undirectable}[1]{\mathbb{L}_{#1}}
\newcommand{\Xzero}[1]{X_{#1}}

% Noun macros

\newcommand{\Baire}{Baire\xspace}
\newcommand{\Borel}{Bor\-el\xspace}
\newcommand{\Hausdorff}{Haus\-dorff\xspace}
\newcommand{\Kechris}{Kech\-ris\xspace}
\newcommand{\Polish}{Po\-lish\xspace}
\newcommand{\Solecki}{So\-lec\-ki\xspace}
\newcommand{\Todorcevic}{To\-dor\-\v{c}ev\-ic\xspace}

% Proof environment
\newenvironment{theoremproof}{
	\renewcommand{\qedsymbol}{$\Box$}
	\begin{proof}
	}{\end{proof}}

% Theorem environments
\newtheorem{theorem}{Theorem}[section]
\newtheorem{lemma}[theorem]{Lemma}
\newtheorem{Claim}[theorem]{Claim}
\newtheorem{claim}[theorem]{Claim}
\newtheorem{subclaim}[theorem]{Subclaim}
\newtheorem{corollary}[theorem]{Corollary}
\newtheorem{proposition}[theorem]{Proposition}
\newtheorem{conjecture}[theorem]{Conjecture}
\newtheorem{question}[theorem]{Question}
\newtheorem{problem}[theorem]{Problem}
\newtheorem{fact}[theorem]{Fact}
\newtheorem{ax}{Axiom}


\theoremstyle{definition}

\newtheorem{definition}[theorem]{Definition}
\newtheorem{example}[theorem]{Example}
\newtheorem{remark}[theorem]{Remark}

\numberwithin{equation}{section}

% More mathematical macros
%\newcommand{\R}{\mathbb{R}}
\newcommand{\Q}{\mathbb{Q}}
%\newcommand{\N}{\mathbb{N}}
\newcommand{\Z}{\mathbb{Z}}
\newcommand{\rank}{\qopname\relax o{rk}}
\newcommand{\dlength}{\qopname\relax o{dilength}}
\newcommand{\length}{\qopname\relax o{length}}
\newcommand{\iB}{\mathcal{B}}
\newcommand{\iF}{\mathcal{F}}
\newcommand{\iU}{\mathcal{U}}
\newcommand{\iV}{\mathcal{V}}
\newcommand{\al}{\alpha}
\newcommand{\la}{\lambda}
\newcommand{\om}{\mathbb{N}}
\newcommand{\omm}{[\mathbb{N}]^\mathbb{N}}
\newcommand{\oom}{\mathbb{N}^\mathbb{N}}
\newcommand{\si}{\sigma}
\newcommand{\eps}{\varepsilon}
\DeclareMathOperator{\maxdeg}{maxdeg}
\newcommand{\RR}{\mathbb{R}}
\newcommand{\bd}{\begin{definition}}
	\newcommand{\ed}{\end{definition}}
\newcommand{\solc}{\mathop{sc}}
\newcommand{\lab}{\label}
\DeclareMathOperator{\graph}{graph}
\DeclareMathOperator{\dom}{dom}
\DeclareMathOperator{\proj}{proj}
\DeclareMathOperator{\ran}{ran}
\DeclareMathOperator{\imm}{im}
\DeclareMathOperator{\rd}{rd}
\DeclareMathOperator{\supp}{supp}
\DeclareMathOperator{\mgs}{mgs}
\DeclareMathOperator{\dist}{dist}
\DeclareMathOperator{\didistance}{didist}
\newcommand{\mc}{\mathcal}
\newcommand{\mb}{\mathbf}
\newcommand{\bs}{\mathbf{\Sigma}^1_1}
\newcommand{\bbg}{\mathbf{\Gamma}}
\newcommand{\bbo}{\mathbf{\Delta}^1_1}
\newcommand{\bp}{\mathbf{\Pi}^1_1}
\newcommand{\ls}{\Sigma^1_1}
\newcommand{\lbo}{\Delta^1_1}
\newcommand{\lp}{\Pi^1_1}
\newcommand{\ve}{\varepsilon}
\newcommand{\mf}{\mathfrak}
\newcommand{\nr}{\lnot R}
\newcommand{\id}{id}
\newcommand{\distance}[3]{\ifthenelse{\isempty{#3}}{\dist(#1,#2)}{\dist^{#3}(#1,#2)}}
\newcommand{\didist}[3]{\ifthenelse{\isempty{#3}}{\didistance(#1,#2)}{\didistance^{#3}(#1,#2)}}
\newcommand{\digraph}[3]{\ifthenelse{\equal{#1}{b}}{\mathbb{#2}_{#3}}
	{{#2}_{#3}}}
\newcommand{\linegraph}[3]{\ifthenelse{\equal{#1}{b}}{\mathbb{#2}_{#3}}
	{#2_{#3}}}
\newcommand{\specialvertex}[2]{s_{#1}}
\newcommand{\underlyingspace}[3]{\ifthenelse{\equal{#1}{b}}{\mathbb{#2}_{#3}}
	{#2_{#3}}}
\newcommand{\distanceset}[2]{\ifthenelse{\isempty{#2}}{D(#1)}{D^{#2}(#1)}}
\newcommand{\Borelchromatic}[1]{\chi_B(#1)}
\newcommand{\restrcitedto}[2]{#1|#2}
\newcommand{\connecteder}[1]{E_{#1}}
%\newcommand{\dist}{dist}
%\newcommand{\length}{length}
%\newcommand{\dlength}{dilength}
\newcommand{\concatt}{%
	\mathbin{\raisebox{1ex}{\scalebox{.7}{$\frown$}}}%
}
\newcommand{\ovl}{\overline}
\newcommand{\pred}{pred}
\newcommand{\directed}[1]{\mathbb{L}_{#1}}
\newcommand{\Hzerodirected}[2]{\mathbb{H}_{#1,#2}}
\newcommand{\aStage}[1]{n^{#1}}
\newcommand{\aCode}[1]{f^{#1}}
\newcommand{\aCodeB}[1]{c^{#1}}
\newcommand{\aCodeD}[1]{d^{#1}}
\newcommand{\aMapV}[1]{\phi^{#1}}
\newcommand{\aMapE}[2]{\psi^{#1}_{#2}}
\newcommand{\aBoundStage}[1]{\aCodeB{#1},\aStage{#1}}
\newcommand{\aDifferent}[2]{\aCodeB{#2},\aStage{#1},\aStage{#2}}
\newcommand{\aCodeStage}[1]{\aCode{#1},\aStage{#1}}
\newcommand{\Lias}[3]{((#2_{i})_{i<#1},(#3_{i})_{0<i<#1})}
\newcommand{\aLias}[2]{#1^{#2}}
\newcommand{\imageConf}[2]{A(#1,#2)}
\newcommand{\BorelInvariant}[2]{B(#1,#2)}
\newcommand{\twocoloring}[2]{c_{#1,#2}}
\newcommand{\oddpair}[1]{odd $#1$-pair}
\newcommand{\lgth}{\length}
%\newcommand{\N}{\mathbb{N}}
\newcommand{\Ramseyspace}{[\N]^{\N}}
%\newcommand{\dist}{dist}
\newcommand{\Shiftmap}{\mathcal{S}}
\newcommand{\Shiftgraph}{\mathcal{G}_\Shiftmap}
\newcommand{\Shiftgraphmodified}{\mathcal{G}'_{\Shiftmap}}
\newcommand{\Homomorphism}[2]{Hom(#1,#2)}
\newcommand{\Homomorphisma}[2]{Hom^*(#1,#2)}
\newcommand{\Tree}[1]{T_{#1}}
\newcommand{\Homothree}[1]{Hom(\Tree{3},#1)}
\newcommand{\Homothreea}[1]{Hom^*(\Tree{3},#1)}
\newcommand{\Deltaabsolute}[1]{\chi_{\Delta^1_2ab}(#1)}
\newcommand{\Bairemeasurablechromatic}[1]{\chi_{BM}(#1)}
\newcommand{\Measurablechromatic}[1]{\chi_{\mu}(#1)}
\newcommand{\Ztwocube}{\mathbb{Z}_2*\mathbb{Z}_2*\mathbb{Z}_2}
\newcommand{\SG}{\mathcal{G}_\mathcal{S}}



\newcommand{\rstr}{\restriction}
\newcommand{\vfi}{\varphi}
\newcommand{\sse}{\subseteq}
\newcommand{\est}{\emptyset}


\newcommand{\treegraph}[1]{Tree(#1)}
\newcommand{\fG}{\mathcal{G}}
\newcommand{\fA}{\mathcal{A}}
\newcommand{\fB}{\mathcal{B}}
\newcommand{\fC}{\mathcal{C}}
\newcommand{\fD}{\mathcal{D}}
\newcommand{\fE}{\mathcal{E}}
\newcommand{\Epsilon}{\fE}
\newcommand{\fF}{\mathcal{F}}
%\newcommand{\fG}{\mathcal{G}}
\newcommand{\fH}{\mathcal{H}}
\newcommand{\fI}{\mathcal{I}}
\newcommand{\fJ}{\mathcal{J}}
\newcommand{\fK}{\mathcal{K}}
\newcommand{\fL}{\mathcal{L}}
\newcommand{\fM}{\mathcal{M}}
\newcommand{\fN}{\mathcal{N}}
\newcommand{\fO}{\mathcal{O}}
\newcommand{\fP}{\mathcal{P}}
\newcommand{\fQ}{\mathcal{Q}}
\newcommand{\fR}{\mathcal{R}}
\newcommand{\fS}{\mathcal{S}}
\newcommand{\fT}{\mathcal{T}}
\newcommand{\fU}{\mathcal{U}}
\newcommand{\fV}{\mathcal{V}}
\newcommand{\fW}{\mathcal{W}}
\newcommand{\fX}{\mathcal{X}}
\newcommand{\fY}{\mathcal{Y}}
\newcommand{\fZ}{\mathcal{Z}}

% Tikz settings

\crefname{originaltheorem}{original theorem}{original theorems}
\Crefname{originaltheorem}{Original theorem}{Original theorems}
\makeatletter
\newlabel{t:separation2}{{1.4}{source}{}{sourcepointer-1}{}}
\newlabel{t:separation2@cref}{{[originaltheorem][0][]1.4}{[1][0][]source}{}{}{}}
\makeatother
\title{Every countable Borel equivalence relation is hyperfinite on a comeagre invariant set}
\author{Balazs Bursics, Zoltan Vidnyanszky}
\date{arxiv:2412.01315v2; Creative Commons Attribution 4.0 International}
\begin{document}
\maketitle
\noindent\textit{Editorial note.} The counted statement is the author's environment at source0.tex lines 664-666 and its proof environment opens at line 668; the pairing is by adjacency and the retained proof range is the whole environment.\par\medskip
\noindent\textit{Editorial note.} Preamble: the author's own package lines and macro block; the drawing and plotting packages of the source are not retained because the retained material does not use them.\par\medskip
\noindent\textit{Editorial note.} Retained-licence form: this source's licence evidence carries no acceptance record and none may be written on anyone's behalf, so the card points at a task-directory evidence copy in the team-wide four-key reuse form; no status key and no acceptance record are written.\par\medskip
\noindent\textit{Editorial note.} Metadata provenance: source\_version is the value recorded in the source's licence evidence (2412.01315v2); the author list and title are transcribed from the authors\_index and title keys of the source record of the pinned copy.\par\medskip
\noindent\textit{Editorial note.} No label-number dump exists for this source and no counter values are imposed, so identity is the author label, or the pinned source position where the author gave none.\par\medskip
\noindent\textit{Editorial note.} r002 note: r001 failed because the author macro block uses newcommandx from the xargs package (and ifthenelse and textcolor), none of which the retained package spans carried; r002 adds those three packages to extra\_packages, as baseline packages of the author, and changes nothing else.\par\medskip
\noindent\textit{Editorial note.} r003 note: the author macro block defines nine macros with newcommandx (source0.tex lines 79-83 concatenation, 89-91 functions, 99-103 Hzero, 104-108 intersection, 118-122 product, 126-130 sequence, 132-136 set, 137-139 sets, 147-151 union), which the xargs package provides; the host TeX Live does not ship xargs.sty, so those definition lines are removed from the preamble span instead of the package being added, and xargs is removed from extra\_packages. A transitive check of every macro referenced by the retained statement and proof spans (operations/agent-authoring/lit/macroclosure.py) shows that no retained span uses any of those nine macros, directly or through another author macro body, so nothing the cards print depends on them. ifthen and xcolor are dropped as well because the author's own xifthen package already supplies ifthenelse and the retained spans expand no textcolor. The nine macro bodies and every other line of the block are retained verbatim.\par\medskip
\noindent\textit{Editorial note.} r003 bibliography note: the retained proof cites KST, and the pinned source directory ships the author bibliography ref.bib (392 entries); the bibliography backend is bibtex with the source own ref.bib attached verbatim as bibliography\_paths. The cited key KST was found in that file at its article entry (entry begins on line 1924) and no entry was edited, added or rewritten. This also exercises the bibtex route on this lane.\par\medskip
\noindent\textit{Editorial note.} r004 note: r003 failed on the host with `LaTeX Error: Environment lemma undefined` because the card retained the author's package lines and macro block but not the block that declares the theorem environments. That block is the author's own text at source0.tex lines 176-204 (the theoremproof environment, the theorem/lemma/claim/subclaim/corollary/proposition/conjecture/question/problem/fact/ax declarations, the definition theoremstyle with its definition/example/remark declarations, and numberwithin for equations) and r004 appends it as one more preamble span, verbatim. Nothing else changed.\par\medskip
\noindent\textit{Editorial note.} r005 note: the retained proof cites t:separation2, which the card does not contain; per the authoring manual 8.1 it is declared in external\_references as a non-load-bearing navigation pointer (auto-generated newlabel stub) rather than pulled in as a context span, because pulling it in would consume a number of the retained sequence. Numbers are inferred from the source own section-numbered counter and are pointers only.\par\medskip
\noindent\textit{Editorial note.} r006 note: r005 passed the theorem declarations and then failed with `Undefined control sequence \textbackslash{}maxdeg`, an operator name the author defines at source0.tex line 225 inside the preamble tail (lines 205-349). The r004 span 176-204 is therefore extended to 176-349, which is the author's own text up to the line before \textbackslash{}begin\{document\} (line 350); every retained line is verbatim and no macro is invented.\par\medskip
\noindent\textit{Editorial note.} r007 note: the card-local licence evidence copy carried neither source\_id nor source\_version, so the packaged licenses/evidence.json could not be compared with this task snapshot (missing key against a value) and the publisher deferred the card as an sv mismatch. The evidence file now carries source\_id and source\_version written from this card by operations/agent-authoring/lit/sv\_guard.py, the single writer for those two fields, so the two sides are equal byte for byte (arxiv:2412.01315v2 / 2412.01315v2, the value the source own LICENSE\_EVIDENCE.json records). The four-key reuse form is unchanged, no status key is written and no acceptance record is written. Nothing else in the card changed; the revision is bumped only so the package is rebuilt with the aligned evidence.\par\medskip
\noindent\textit{Editorial scope.} One original result of the pinned author manuscript, counted once: hyperfiniteness of a countable Borel equivalence relation on a comeagre invariant set (unlabelled theorem at source0.tex lines 664-666).\par\medskip
\noindent\textit{Allowed standard prerequisites (editorial).}\begin{itemize}
\item Countable Borel equivalence relation. A Borel equivalence relation all of whose classes are countable, expressible as a countable union of graphs of Borel involutions.. It is the relation whose restriction is shown to be hyperfinite.
\item Hyperfinite relation. An equivalence relation is hyperfinite when it is the increasing union of finite Borel equivalence relations.. It is the property concluded on the restricted relation.
\item Comeagre invariant set. A comeagre Borel set that is a union of classes, so that the restriction of the relation to it is still an equivalence relation.. It is the set produced by the construction.
\end{itemize}\medskip
\ifdefined\hypertarget\hypertarget{sourcepointer-1}{}\fi
\noindent\textit{Source reference.} Original theorem 1.4: sources/hyperfinite-ramsey-spaces-2412.01315v2/source0.tex, source lines 426--426. Original result of the same manuscript invoked by the retained proof at source0.tex line 426 as a navigation pointer (kind theorem, original number 1.4). Nothing load-bearing is replaced by it: the counted statement's own author proof is retained line for line and complete. The number was inferred by counting, inside the enclosing section, the theorem-class environments that share the same section-numbered counter declared by the source's own newtheorem lines (operations/agent-authoring/lit/envnum.py); this source carries no label-number dump, so the number must be read as a pointer only.\par\medskip
	\begin{theorem}
		Let $E$ be a CBER on the space $X$. Then there is a comeager invariant Borel set $C$ such that $\rstr{E}{C}$ is hyperfinite.
	\end{theorem}

	\begin{proof}
		As all classes of $E$ are countable, there are Borel involutions $\vfi_n$ on $X$ such that $E=\bigcup_{n\in\N}\graph(\vfi_n)$. Let $G_n=\bigcup_{i\in n}\graph(\vfi_i)$, and $f:\N\to \N$ be as in Theorem \ref{t:separation2}. We will construct Borel sets $B_n \subseteq X$ so that every $B_n$ is $f(n)$-separated in $G_n$, and there is a comeager invariant Borel set $C$ such that $C\sse\{x:\exists^\infty n \ (x \in B_n)\}$. 
		
		Fix a bijection $g:\N\to\N^2$ with coordinate functions $g_0,g_1$ and enumerate a basis of $X$ as $(U_n)_{n\in\N}.$ Since \[G_n^{f(n)}=\{(x,y):x\neq y \wedge \dist_{G_n}(x,y)\le f(n) \}\]  has bounded degrees there exists a finite Borel vertex coloring $c_n:X\to m$ of $G_n^{f(n)}$ by \cite{KST}. For some $i\in m$ the preimage $c_n^{-1}(i)$ is non-meager in $U_{g_0(n)}.$ Set $B_n=c_n^{-1}(i).$ This way $B_n$ is $f(n)$-separated in $G_n$ and for every $k\in\N$ we have that
		\[C_k=\bigcup_{n\in g_1^{-1}(k)}B_n\]
		is comeager in $X,$ and 
		\[C'=\bigcap_{k\in\N}C_k\sse \{x:\exists^\infty n \ (x \in B_n)\}.\]
		As the saturation of $X\setminus C'$ is meager, $C =C'\setminus\big[X\setminus C'\big]_E$ is a comeager invariant Borel set such that $\rstr{E}{C}$ is hyperfinite by Theorem \ref{t:separation2}.
	\end{proof}

\bibliographystyle{plain}
\bibliography{ref}
\end{document}
